\section{Dataset Construction and Shared Representation}\label{sec:dataset}

\subsection{Goals and construction paths}
The representation is designed first for route-aware ETA learning while remaining useful beyond a single supervised target: it places road topology, vehicles, route context, and time-varying traffic state in one coherent form for controlled experimentation and downstream testing. Because it preserves graph semantics rather than only a final ETA label, it can also support related tasks such as road-level load estimation and area-level analysis.

In the simulation-derived path, construction begins from a controlled microscopic environment built with SUMO and managed through the desktop tool and TraCI~\cite{sumo2011,traci2018}. Planned routes, vehicle identities, network structure, and evolving traffic states are directly observable during execution. The desktop workflow is configurable: simulation behavior, validation of zones, landmarks, and vehicle classes, temporal sampling of snapshots, and optional irregularity such as random untracked trips for robustness experiments.
Figure~\ref{fig:sim_path_ui}(\textbf{a}) shows the simulation-definition stage; Figure~\ref{fig:sim_path_ui}(\textbf{b}) shows the web client used to evaluate predictors on simulation-backed journeys.

% Simulation branch: desktop zone/scenario definition (a) + web ETA evaluation on simulation exports (b).
\begin{figure}[t]
\centering
\newlength{\simPathFigHeight}
\setlength{\simPathFigHeight}{3.35cm}
\begin{minipage}[c]{0.49\linewidth}
\centering
\includegraphics[width=\linewidth,height=\simPathFigHeight,keepaspectratio]{dataset_zones.png}\\[-0.05em]
{\scriptsize (\textbf{a}) Simulation path: zones and scenario setup (desktop)}
\end{minipage}%
\hfill
\begin{minipage}[c]{0.49\linewidth}
\centering
\includegraphics[width=\linewidth,height=\simPathFigHeight,keepaspectratio]{web_sim.png}\\[-0.05em]
{\scriptsize (\textbf{b}) Web client: evaluation on simulation-backed trips}
\end{minipage}
\caption{Simulation construction path and its web evaluation loop. (\textbf{a})~Simulation Definition: SUMO network map with attribute-based zones (A/B/C/H), zone statistics, and configuration before database generation and runs (example project). (\textbf{b})~Real-time ETA testing: model performance summaries, interactive map with active route, and journey log (demonstrator layout; branding omitted in figures).}
\label{fig:sim_path_ui}
\label{fig:web_eval_dashboard}
\end{figure}


The real-trajectory path starts from noisier traces: GPS inaccuracies, temporal inconsistencies, or jumps may break a coherent path interpretation. Conversion is a staged repair-and-alignment workflow: network preparation, import, repair, route inference or matching, slicing of trajectories with major jumps, and filtering of degenerate origin--destination pairs.
We use the Porto taxi trajectory dataset~\cite{porto2015pkdd} as a representative real-world conversion example in our open-source bundle; the toolchain applies to other CSVs and regions.
Figure~\ref{fig:traj_path_ui}(\textbf{a}) shows the desktop conversion interface; Figure~\ref{fig:traj_path_ui}(\textbf{b}) shows web-side evaluation on trajectory-backed runs.
Figure~\ref{fig:route_matching} visualizes route-matching outputs on the map: search scope, candidates, raw GPS, the chosen SUMO route, and projections onto it.

% Trajectory branch: desktop conversion (a) + web evaluation on Porto-aligned exports (b).
\begin{figure}[t]
\centering
\newlength{\trajPathFigHeight}
\setlength{\trajPathFigHeight}{3.35cm}
\begin{minipage}[c]{0.49\linewidth}
\centering
\includegraphics[width=\linewidth,height=\trajPathFigHeight,keepaspectratio]{dataset_convertion_main_with_map.png}\\[-0.05em]
{\scriptsize (\textbf{a}) Trajectory path: conversion (desktop)}
\end{minipage}%
\hfill
\begin{minipage}[c]{0.49\linewidth}
\centering
\includegraphics[width=\linewidth,height=\trajPathFigHeight,keepaspectratio]{porto_web.png}\\[-0.05em]
{\scriptsize (\textbf{b}) Web client: evaluation on trajectory-backed trips}
\end{minipage}
\caption{Trajectory construction path and its web evaluation loop. (\textbf{a})~Trajectory dataset conversion: OSM/SUMO map, network overlay, CSV inputs, and batch conversion controls (example Porto-area project). (\textbf{b})~Same web client as in Figure~\ref{fig:sim_path_ui}(\textbf{b}), here exercised on exports aligned with real trajectories: map with active route, aggregate metrics, and recent journeys (demonstrator layout; branding omitted in figures).}
\label{fig:traj_path_ui}
\end{figure}


% Route matching: bounding-box / search scope (a) and split trajectory repair (b).
\begin{figure}[t]
\centering
\newlength{\routeMatchFigHeight}
\setlength{\routeMatchFigHeight}{3.55cm}
\begin{minipage}[t]{0.49\linewidth}
\centering
\includegraphics[width=\linewidth,height=\routeMatchFigHeight,keepaspectratio]{dataset_conversion_bounding_box_whitebg.png}\\[-0.2em]
{\scriptsize (\textbf{a}) Bounding box and search scope}
\end{minipage}%
\hfill
\begin{minipage}[t]{0.49\linewidth}
\centering
\includegraphics[width=\linewidth,height=\routeMatchFigHeight,keepaspectratio]{dataset_conversion_split_whitebg.png}\\[-0.2em]
{\scriptsize (\textbf{b}) Split trajectory after a GPS outlier}
\end{minipage}
\caption{Route matching on the imported SUMO network (conversion path). Panels~(\textbf{a}) and~(\textbf{b}) both show the \emph{bounding box}, \emph{burgundy} roads and junctions that limit the matcher's search entities, and the thick \emph{dark} final route on the SUMO network. Panel~(\textbf{a}) additionally shows the raw GPS trace as a \emph{light blue} polyline (numbered blue samples) and \emph{green} segments as the top-three candidate edges at sampled locations. Panel~(\textbf{b}) additionally shows \emph{light blue} and \emph{orange} polylines for the \emph{split} trajectory---two segment polylines after cutting at an outlying GPS point (e.g.\ $\sim\!5$\,km from adjacent samples), with each segment matched separately.}
\label{fig:route_matching}
\end{figure}


\subsection{Route matching on the network (abstract procedure)}
Matching is not run on the full regional road graph: first, an axis-aligned bounding box $B$ is built around the GPS polyline $P$ (optionally padded), and $N_{\mathrm{bb}}=(V_{\mathrm{bb}},E_{\mathrm{bb}})$ is the subnetwork of $N$ formed from all edges whose geometry intersects $B$ and their incident nodes.
GPS polylines are then aligned to $N_{\mathrm{bb}}$ in two phases: a \emph{scoring pass} assigns every edge a coarse tier ($1$, $10$, or $1000$) from successive segment--edge comparisons, then a \emph{routing pass} searches for a low-cost path that respects those tiers.
For each polyline segment, edges whose heading differs from the segment by more than $\pi/9$ are discarded; among survivors, edges are ranked by $r_e=\lambda\delta_e+\frac{1}{2}\bigl(d(p_i,e)+d(p_{i+1},e)\bigr)$, where $\delta_e$ is the acute angle gap in radians, $d(p,e)$ is the distance from point $p$ to the centreline polyline of edge $e$, and $\lambda$ fixes the radian-to-metre (map-unit) trade-off.
The first $k$ edges become \emph{primary} candidates (tier~$1$), the next $j$ become \emph{secondary} (tier~$10$); edges that never qualify as a potential match in any segment keep tier~$1000$.
Whenever an edge appears among the top-$k$ candidates for a segment, its tier is set to $1$ outright; secondary ($10$) updates use $\min$ with the current tier so an edge that already has tier $1$ is never raised to $10$ by a later segment.
The routing pass runs A$^\ast$ exactly $k$ times: each run starts from one edge in $\mathrm{Top}_1$ (the $k$ primaries for the first segment) and \emph{terminates} as soon as the search reaches any of the \emph{last} primaries---any edge in $\mathrm{Top}_{m-1}$ for segment $(p_{m-1},p_m)$---which form the goal set. Edge costs are tier-weighted; the heuristic is Euclidean distance to $p_m$. The cheapest route among the $k$ runs is returned.
Algorithms~\ref{alg:route-edge-weights}--\ref{alg:route-astar} state this procedure formally.

% Route matching: abstract scoring + A* as specified for the paper (conceptual layer).

\begin{algorithm}[htbp]
\caption{Edge-weight assignment from a GPS polyline (scoring pass)}
\label{alg:route-edge-weights}
\begin{algorithmic}[1]
\Require Polyline $P=(p_1,\ldots,p_m)$; full road network $N=(V,E)$ with directed edges;
         integers $k,j \ge 1$; angular threshold $\theta_{\max}=\pi/9$; trade-off $\lambda>0$ (same length units as the map, e.g.\ $\lambda=5$).
\Ensure Each edge $e\in E_{\mathrm{bb}}$ has a tier $w(e)\in\{1,10,1000\}$ (lower is more attractive).
\State Build axis-aligned bounding box $B$ containing all points of $P$ (optional margin).
\State Let $N_{\mathrm{bb}}=(V_{\mathrm{bb}},E_{\mathrm{bb}})\subseteq N$ be the subnetwork of all edges in $E$ whose centreline intersects $B$, together with their endpoint nodes in~$V$.
\State Initialise $w(e)\gets 1000$ for all $e\in E_{\mathrm{bb}}$.
\For{each consecutive segment $s_i=(p_i,p_{i+1})$, $i=1,\ldots,m-1$}
    \State $\phi_i \gets \mathrm{atan2}\bigl(p_{i+1}-p_i\bigr)$ \Comment{segment heading in the projected map}
    \State $\mathcal{C}\gets\emptyset$ \Comment{candidate edges for this segment}
    \For{each edge $e\in E_{\mathrm{bb}}$}
        \State $\psi_e \gets$ heading of $e$ from its centreline polyline (first to last shape point)
        \State $\delta_e \gets \min\bigl(|\phi_i-\psi_e|,\,2\pi-|\phi_i-\psi_e|\bigr)$
        \If{$\delta_e \le \theta_{\max}$}
            \State mark $e$ as \emph{potential} for $s_i$; add $e$ to $\mathcal{C}$
            \State $d_i \gets d(p_i,e)$; $d_{i+1} \gets d(p_{i+1},e)$
            \Comment{shortest Euclidean distance from each segment endpoint to the polyline of $e$}
            \State $r_e \gets \lambda\,\delta_e + \tfrac{1}{2}(d_i + d_{i+1})$
            \Comment{$\lambda$ scales the angle gap $\delta_e$ (radians) into map units so it is comparable to the mean endpoint--edge distances.}
        \EndIf
    \EndFor
    \State sort $\mathcal{C}$ by increasing $r_e$
    \State $\mathrm{Top}_i \gets$ first $k$ edges; $\mathrm{Sec}_i \gets$ next $j$ edges (if $|\mathcal{C}|<k{+}j$, take as many as exist)
    \For{$e\in \mathrm{Top}_i$}
        \State $w(e)\gets 1$
    \EndFor
    \For{$e\in \mathrm{Sec}_i$}
        \State $w(e)\gets \min\bigl(w(e),10\bigr)$
    \EndFor
\EndFor
\Comment{Edges that were never marked potential in any segment keep $w(e)=1000$. Each time an edge ranks in $\mathrm{Top}_i$, set $w(e)\gets 1$; secondary ranks use $\min$ with $10$ so a prior tier $1$ is not overwritten by a weaker segment.}
\end{algorithmic}
\end{algorithm}

\begin{algorithm}[htbp]
\caption{Route extraction by multi-start A$^\ast$}
\label{alg:route-astar}
\begin{algorithmic}[1]
\Require Restricted network $N_{\mathrm{bb}}$; tiers $w(e)$ on $E_{\mathrm{bb}}$; sets $\mathrm{Top}_1$ and $\mathrm{Top}_{m-1}$ from Algorithm~\ref{alg:route-edge-weights} (primaries for the first and last polyline segments); last point $p_m$.
\Ensure Shortest-cost route among $k$ multi-start runs.
\State Goal set $G \gets \mathrm{Top}_{m-1}$ \Comment{the $k$ primary edges for the last segment $(p_{m-1},p_m)$}
\For{each start edge $e_{\mathrm{start}}\in \mathrm{Top}_1$}
    \State Run A$^\ast$ from $e_{\mathrm{start}}$ with edge cost $c(e)=w(e)\cdot \mathrm{len}(e)$ and $h(v)=$ Euclidean distance from $v$ to $p_m$;
           \textbf{stop} the run as soon as the search reaches any state on an edge in $G$.
\EndFor
\State \Return the route with minimum total cost among these $k$ runs.
\end{algorithmic}
\end{algorithm}


Finally, GPS samples are projected onto the inferred SUMO route so vehicle position can be aggregated as a time series in the shared representation.

\subsection{Shared representation, dynamics, and use}
Both paths emit dataset instances in the same conceptual form: roads as edges, junctions and vehicles as nodes, with planned or inferred routes so the graph encodes route-conditioned traffic context, not only a traffic snapshot. Exports are feature-rich at node and edge level. Figure~\ref{fig:graph_schema}(\textbf{b}) illustrates static junctions and roads together with vehicle nodes and dynamic edges (vehicle--vehicle and vehicle--junction).

Each snapshot is tied to a simulation step or to a trajectory update instant and bundles the static road/junction substrate with the active vehicle set, its routes, and the current traffic descriptors on touched edges and nodes. That alignment matters for learning: models trained on SUMO exports and models trained on converted trajectories ingest the same logical tensors even though the emission times are regular in one branch and event-driven in the other.

The graph is dynamic in two senses: node and edge attributes evolve with traffic conditions, and structure changes as vehicle nodes move and dynamic relations are created or removed. The stable road graph is supplemented by transient traffic relations reflecting positioning and interaction. Downstream users may take the full dynamic graph or simplified views (vehicles, roads, junctions; optional filtering of dynamic relations).
